\documentclass[10pt,a4paper]{amsart}
\usepackage[margin=19mm]{geometry}
\usepackage{amsmath,amssymb,mathtools}
\usepackage[scr=rsfs,cal=euler]{mathalfa}
\usepackage{xfrac}
\usepackage[unicode,hidelinks,bookmarks=false]{hyperref}
\newcommand{\ie}{i.e.\ }
\newcommand{\cf}{cf.\ }
\newcommand{\resp}{respectively\ }
\newcommand{\Subsection}{\subsection}
\providecommand{\nobreakdash}{\nobreak\hbox{-}}
\setlength{\emergencystretch}{2em}
\multlinegap0pt
\usepackage{mathtools}
\usepackage{bm}
\usepackage{algorithm}\usepackage{algpseudocode}
\newtheorem{thm}{Theorem}
\newcommand{\mb}{\boldsymbol}
\title{Deep Adaptive Dimension Reduction for Bayesian Inference in Inverse Problems}
\author{Yueyang Wang, Xili Wang, Kejun Tang, Xiaoliang Wan, Tao Zhou,, Chao Yang}
\date{Source: arXiv:2605.29373v2 (CC BY 4.0)}
\begin{document}
\maketitle

\noindent\emph{Source and licence.} Deep Adaptive Dimension Reduction for Bayesian Inference in Inverse Problems. Version arXiv:2605.29373v2 (CC BY 4.0), distributed by arXiv under the Creative Commons Attribution 4.0 International licence, http://creativecommons.org/licenses/by/4.0/, as recorded in the arXiv licence metadata for this exact version and shown on the abstract page https://arxiv.org/abs/2605.29373v2. Full licence notice: \texttt{licenses/arxiv-2605.29373-CC-BY-4.0.txt}.\par\smallskip

\noindent\textit{Editorial scope.} Counted unit: the article's thm thm:vf\_elbo\_extended (source0.tex lines 401--418). The article's complete original proof of it is delivered with it (source0.tex lines 419--519, one \texttt{proof} environment of 101 lines). No mathematics is written, repaired or re-derived: the statement and the proof are the article's own lines, in the article's own notation, and no retained line is substituted or re-flowed.  No further source-local context is needed: every cross-reference of every retained line resolves inside the two delivered spans.  Nothing was omitted from any retained span, so the package carries no omission record.  No retained line contains a citation, so the wrapper carries no bibliography and no bibliography entry.  No retained line contains a figure environment, an image-inclusion command or drawing code, so no image is delivered and the package declares no asset dependency.  Editorial formatting only, in addition: an AMS \texttt{amsart} preamble that restates the article's own declarations and macros used by the retained lines, namely the environments \texttt{thm} on separate counters, together with the standard packages those lines need.  The retained environments print in this document as Theorem 1 in order of appearance, whereas the article prints the same items with the numbering of its own full document.  The complete native source of the article as distributed in the arXiv e-print for this exact version is bundled unchanged as \texttt{sources/201605/source0.tex}.
\begin{thm}
\label{thm:vf_elbo_extended}
Suppose the decoder $p_{\mb{x}|\mb{z},\mb{\theta}}$ in VAE is parameterized by nonlinear neural networks. If
\begin{enumerate}
\item[(a)] the family of flow priors $\{p_{\mb{z},\mb{\beta}}\}_{\mb{\beta}}$ induced by the bijection $f_{\mathrm{pr},\boldsymbol{\beta}}$ can approximate any $k$-dimensional probability density function to arbitrary precision;
\item[(b)] the family of conditional normalizing flow encoders $\{q_{\mb{z}|\mb{x},\boldsymbol{\alpha}}\}_{\boldsymbol{\alpha}}$ can approximate any conditional distribution to arbitrary precision;
\end{enumerate}
suppose $(\mb{\phi}^*,\mb{\theta}^*)$ are the optimal parameters for standard VAE:
\begin{equation}
(\mb{\phi}^*,\mb{\theta}^*) = \mathop{\arg\max}_{\mb{\phi},\, \mb{\theta}} \; \mathbb{E}_{p_{\mb{x}}}[L_{\mb{\phi},\mb{\theta}}(\mb{x})],
\end{equation}
then the Variational Flow (VF) model achieves a strictly higher Evidence Lower Bound (ELBO) than standard VAE if either of the following conditions holds:
\begin{enumerate}
\item[(i)] the aggregated posterior of the optimal encoder in the standard VAE does not equal the prior, i.e., $q_{\mb{z},\boldsymbol{\phi}^*}= \int q_{\mb{z}|\mb{x},\mb{\phi}^*} p_{\mb{x}}d\mb{x} \neq p_{\mb{z}}$;
\item[(ii)] there exists a set of positive measure with respect to $p_{\mb{x}}$ such that the model posterior $p_{\mb{z}|\mb{x},\boldsymbol{\theta}^*}=p_{\mb{x}|\mb{z},\mb{\theta}^*}p_{\mb{z}}/ \int p_{\mb{x}|\mb{z},\mb{\theta}^*}p_{\mb{z}}d\mb{z}$ of VAE is not a diagonal Gaussian distribution.
\end{enumerate}
Furthermore, when conditions (i) and (ii) hold simultaneously, the flow prior and the conditional normalizing flow encoder each provide a strictly positive improvement to the ELBO.
\end{thm}

\begin{proof}

Let 
\begin{equation}
E_0 \coloneqq \mathbb{E}_{p_{\mb{x}}}[L_{\mb{\phi}^*,\mb{\theta}^*}(\mb{x})]
= -\mathbb{E}_{p_{\mb{x}}}[D_{\mathrm{KL}}(q_{\mb{z}|\mb{x},\mb{\phi}^*} \| p_{\mb{z}})]
+ \mathbb{E}_{p_{\mb{x}}\, q_{\mb{z}|\mb{x},\mb{\phi}^*}}[\log p_{\mb{x}|\mb{z},\mb{\theta}^*}].
\end{equation}

\paragraph{Generalize the prior}
Fixing the encoder $q_{\mb{z}|\mb{x},\mb{\phi}^*}$ and decoder $p_{\mb{x}|\mb{z},\mb{\theta}^*}$ and replacing the prior $p_{\mb{z}}$ with the normalizing flow prior $p_{\mb{z},\mb{\beta}}$, the ELBO becomes
\begin{equation}
\mathbb{E}_{p_{\mb{x}}}[L_{\mb{\phi}^*,\mb{\theta}^*,\mb{\beta}}(\mb{x})]
= -\mathbb{E}_{p_{\mb{x}}}[D_{\mathrm{KL}}(q_{\mb{z}|\mb{x},\mb{\phi}^*} \| p_{\mb{z},\mb{\beta}})]
+ \mathbb{E}_{p_{\mb{x}}\, q_{\mb{z}|\mb{x},\mb{\phi}^*}}[\log p_{\mb{x}|\mb{z},\mb{\theta}^*}].
\end{equation}
Subtracting $E_0$ yields
\begin{equation}
\mathbb{E}_{p_{\mb{x}}}[L_{\mb{\phi}^*,\mb{\theta}^*,\mb{\beta}}(\mb{x})] - E_0
= \mathbb{E}_{p_{\mb{x}}}[D_{\mathrm{KL}}(q_{\mb{z}|\mb{x},\mb{\phi}^*} \| p_{\mb{z}})]
- \mathbb{E}_{p_{\mb{x}}}[D_{\mathrm{KL}}(q_{\mb{z}|\mb{x},\mb{\phi}^*} \| p_{\mb{z},\mb{\beta}})].
\label{eq:diff_prior}
\end{equation}
Simplifying the right-hand side of the above equation, we note that
\begin{align}
\mathbb{E}_{p_{\mb{x}}}[D_{\mathrm{KL}}(q_{\mb{z}|\mb{x},\mb{\phi}^*} \| p)]
&= \int q_{\mb{z}|\mb{x},\mb{\phi}^*}(\mb{z}|\mb{x})\, p_{\mb{x}}(\mb{x}) \log \frac{q_{\mb{z}|\mb{x},\mb{\phi}^*}(\mb{z}|\mb{x})}{p(\mb{z})}\, d\mb{z}\, d\mb{x} \nonumber\\
&= \int q_{\mb{z}|\mb{x},\mb{\phi}^*}\, p_{\mb{x}} \log q_{\mb{z}|\mb{x},\mb{\phi}^*}\, d\mb{z}\, d\mb{x}
- \int q_{\mb{z},\mb{\phi}^*}(\mb{z}) \log p(\mb{z})\, d\mb{z},
\end{align}
where $q_{\mb{z},\mb{\phi}^*}(\mb{z}) = \int q_{\mb{z}|\mb{x},\mb{\phi}^*}(\mb{z}|\mb{x})\, p_{\mb{x}}(\mb{x})\, d\mb{x}$ is the aggregated posterior. Thus, \eqref{eq:diff_prior} simplifies to
\begin{equation}
\mathbb{E}_{p_{\mb{x}}}[L_{\mb{\phi}^*,\mb{\theta}^*,\mb{\beta}}(\mb{x})] - E_0
= D_{\mathrm{KL}}(q_{\mb{z},\mb{\phi}^*} \| p_{\mb{z}}) - D_{\mathrm{KL}}(q_{\mb{z},\mb{\phi}^*} \| p_{\mb{z},\mb{\beta}}).
\label{eq:prior_improvement}
\end{equation}
By assumption (a), the flow prior $p_{\mb{z},\mb{\beta}}$ can approximate $q_{\mb{z},\mb{\phi}^*}$ to arbitrary precision. If condition (i) holds, i.e., $q_{\mb{z},\mb{\phi}^*} \neq p_{\mb{z}}$, then $D_{\mathrm{KL}}(q_{\mb{z},\mb{\phi}^*} \| p_{\mb{z}}) > 0$. By choosing $\tilde{\mb{\beta}}$ such that $D_{\mathrm{KL}}(q_{\mb{z},\mb{\phi}^*} \| p_{\mb{z},\tilde{\mb{\beta}}}) < D_{\mathrm{KL}}(q_{\mb{z},\mb{\phi}^*} \| p_{\mb{z}})$, we have
\begin{equation}
E_1 \coloneqq \mathbb{E}_{p_{\mb{x}}}[L_{\mb{\phi}^*,\mb{\theta}^*,\tilde{\mb{\beta}}}(\mb{x})] > E_0.
\label{eq:E1_gt_E0}
\end{equation}

\paragraph{Generalize the encoder} 
Fixing the decoder $p_{\mb{x}|\mb{z},\mb{\theta}^*}$ and the prior $p_{\mb{z},\tilde{\mb{\beta}}}$, we extend the encoder from a diagonal Gaussian to a conditional normalizing flow $q_{\mb{z}|\mb{x},\boldsymbol{\alpha}}$. Utilizing the following decomposition of the ELBO: 
\begin{align}
\mathbb{E}_{p_{\mb{x}}}[L_{\mb{\alpha},\mb{\theta}^*,\tilde{\mb{\beta}}}(\mb{x})]
&= \mathbb{E}_{p_{\mb{x}}}\left[\int q_{\mb{z}|\mb{x},\mb{\alpha}}\left[\log\frac{p_{\mb{x}|\mb{z},\mb{\theta}^*}p_{\mb{z},\tilde{\mb{\beta}}}}{q_{\mb{z}|\mb{x},\mb{\alpha}}}\right]d\mb{z}\right] \nonumber\\
&= \mathbb{E}_{p_{\mb{x}}}\left[\int q_{\mb{z}|\mb{x},\mb{\alpha}}\left[\log\frac{p_{\mb{x},\mb{\theta}^*,\tilde{\mb{\beta}}}p_{\mb{z}|\mb{x},\mb{\theta}^*,\tilde{\mb{\beta}}}}{q_{\mb{z}|\mb{x},\mb{\alpha}}}\right]d\mb{z}\right] \nonumber\\
&=\mathbb{E}_{p_{\mb{x}}}[\log p_{\mb{x},\mb{\theta}^*,\tilde{\mb{\beta}}}(\mb{x})]
- \mathbb{E}_{p_{\mb{x}}}[D_{\mathrm{KL}}(q_{\mb{z}|\mb{x},\mb{\alpha}} \| p_{\mb{z}|\mb{x},\mb{\theta}^*,\tilde{\mb{\beta}}})],\label{eq:elbo_decomp}
\end{align}
where
\begin{equation}
p_{\mb{x},\mb{\theta}^*,\tilde{\mb{\beta}}}(\mb{x}) = \int p_{\mb{x}|\mb{z},\mb{\theta}^*}(\mb{x}|\mb{z})\, p_{\mb{z},\tilde{\mb{\beta}}}(\mb{z})\, d\mb{z}
\end{equation}
is the model evidence (independent of the encoder parameters), and
\begin{equation}
p_{\mb{z}|\mb{x},\mb{\theta}^*,\tilde{\mb{\beta}}}(\mb{z}|\mb{x})
= \frac{p_{\mb{x}|\mb{z},\mb{\theta}^*}(\mb{x}|\mb{z})\, p_{\mb{z},\tilde{\mb{\beta}}}(\mb{z})}{p_{\mb{x},\mb{\theta}^*,\tilde{\mb{\beta}}}(\mb{x})}
\end{equation}
is the true model posterior under the decoder $p_{\mb{x}|\mb{z},\mb{\theta}^*}$ and the prior $p_{\mb{z},\tilde{\mb{\beta}}}$.
From the decomposition~\eqref{eq:elbo_decomp}, maximizing the ELBO with respect to the encoder is equivalent to minimizing $\mathbb{E}_{p_{\mb{x}}}[D_{\mathrm{KL}}(q_{\mb{z}|\mb{x},\mb{\alpha}} \| p_{\mb{z}|\mb{x},\mb{\theta}^*,\tilde{\mb{\beta}}})]$. In prior generalization, the encoder remains a diagonal Gaussian $q_{\mb{z}|\mb{x},\mb{\phi}^*}$, and the corresponding ELBO is
\begin{equation}
E_1 = \mathbb{E}_{p_{\mb{x}}}[\log p_{\mb{x},\mb{\theta}^*,\tilde{\mb{\beta}}}(\mb{x})]
- \mathbb{E}_{p_{\mb{x}}}[D_{\mathrm{KL}}(q_{\mb{z}|\mb{x},\mb{\phi}^*} \| p_{\mb{z}|\mb{x},\mb{\theta}^*,\tilde{\mb{\beta}}})].
\end{equation}
Since the decoder mean $\mb{\mu}_{\mathrm{de},\mb{\theta}^*}(\mb{z})$ is a nonlinear function of $\mb{z}$, the log model posterior
$$
\log p_{\mb{z}|\mb{x},\mb{\theta}^*,\tilde{\mb{\beta}}}(\mb{z}|\mb{x})
= \log p_{\mb{x}|\mb{z},\mb{\theta}^*}(\mb{x}|\mb{z}) + \log p_{\mb{z},\tilde{\mb{\beta}}}(\mb{z}) - \log p_{\mb{x},\mb{\theta}^*,\tilde{\mb{\beta}}}(\mb{x})
$$
is generally not a quadratic function of $\mb{z}$; thus, $p_{\mb{z}|\mb{x},\mb{\theta}^*,\tilde{\mb{\beta}}}$ is generally not a diagonal Gaussian distribution. If condition (ii) holds, the diagonal Gaussian encoder cannot accurately match the model posterior, i.e.,
\begin{equation}
\mathbb{E}_{p_{\mb{x}}}[D_{\mathrm{KL}}(q_{\mb{z}|\mb{x},\mb{\phi}^*} \| p_{\mb{z}|\mb{x},\mb{\theta}^*,\tilde{\mb{\beta}}})] > 0.
\label{eq:gauss_gap}
\end{equation}
By assumption (b), the conditional normalizing flow encoder $q_{\mb{z}|\mb{x},\mb{\alpha}}$ can approximate any conditional distribution to arbitrary precision. Therefore, there exists $\tilde{\mb{\alpha}}$ such that
\begin{equation}
\mathbb{E}_{p_{\mb{x}}}[D_{\mathrm{KL}}(q_{\mb{z}|\mb{x},\tilde{\mb{\alpha}}} \| p_{\mb{z}|\mb{x},\mb{\theta}^*,\tilde{\mb{\beta}}})]
< \mathbb{E}_{p_{\mb{x}}}[D_{\mathrm{KL}}(q_{\mb{z}|\mb{x},\mb{\phi}^*} \| p_{\mb{z}|\mb{x},\mb{\theta}^*,\tilde{\mb{\beta}}})].
\label{eq:encoder_improvement}
\end{equation}
From \eqref{eq:elbo_decomp} and \eqref{eq:encoder_improvement}, we have
\begin{equation}
E_2 \coloneqq \mathbb{E}_{p_{\mb{x}}}[L_{\tilde{\mb{\alpha}},\mb{\theta}^*,\tilde{\mb{\beta}}}(\mb{x})] > E_1.
\label{eq:E2_gt_E1}
\end{equation}

\paragraph{Combining the two steps}
From \eqref{eq:E1_gt_E0} and \eqref{eq:E2_gt_E1}, when both conditions (i) and (ii) hold,
\begin{equation}
E_2 > E_1 > E_0.
\end{equation}
Since the optimal ELBO of the VF model is no lower than the ELBO value at any specific parameters, we have
\begin{equation}
\sup_{\mb{\alpha},\, \mb{\theta},\, \mb{\beta}} \mathbb{E}_{p_{\mb{x}}}[L_{\mb{\alpha},\mb{\theta},\mb{\beta}}(\mb{x})]
\geq E_2 > E_0
= \sup_{\mb{\phi},\, \mb{\theta}} \mathbb{E}_{p_{\mb{x}}}[L_{\mb{\phi},\mb{\theta}}(\mb{x})].
\end{equation}
When only condition (ii) holds, one can directly choose $\tilde{\mb{\beta}}$ in encoder generalization step to correspond to the identity mapping (i.e., the prior remains $p_{\mb{z}}$), in which case the encoder improvement still provides a strict improvement. In summary, if either condition (i) or (ii) holds, a strict improvement in the ELBO is guaranteed.
\end{proof}

\end{document}
